\documentclass[10pt,a4paper]{amsart}
\usepackage[margin=19mm]{geometry}
\usepackage{amsmath,amssymb,mathtools}
\usepackage[scr=rsfs,cal=euler]{mathalfa}
\usepackage{xfrac}
\usepackage[unicode,hidelinks,bookmarks=false]{hyperref}
\newcommand{\ie}{i.e.\ }
\newcommand{\cf}{cf.\ }
\newcommand{\resp}{respectively\ }
\newcommand{\Subsection}{\subsection}
\providecommand{\nobreakdash}{\nobreak\hbox{-}}
\setlength{\emergencystretch}{2em}
\multlinegap0pt
\usepackage{bm}
\usepackage{bbm}
\usepackage{dsfont}
\usepackage{xspace}
\usepackage{cancel}
\usepackage{mathtools}
\usepackage{amsthm}
\makeatletter
\@ifundefined{thm}{\newtheorem{thm}{Thm}}{}
\@ifundefined{lem}{\newtheorem{lem}[thm]{Lemma}}{}
\makeatother
\newcommand{\RR}{\mathbb{R}}
\newcommand{\p}{{\bf p}}
\title[Expansion of trivariate polynomials using proximity]{Expansion of trivariate polynomials using proximity}
\author{Orit E. Raz}
\date{Source: arXiv:2510.12191v1 (CC BY 4.0)}
\begin{document}
\maketitle

\noindent\emph{Source and licence.} Expansion of trivariate polynomials using proximity. Version arXiv:2510.12191v1 (CC BY 4.0), distributed under the Creative Commons Attribution 4.0 International licence, as recorded in the arXiv licence metadata for this exact version and shown on the abstract page https://arxiv.org/abs/2510.12191v1. Full rights notice: licenses/arxiv-2510.12191-CC-BY-4.0.txt.\par\smallskip

\noindent\textit{Editorial scope.} Counted unit: the Lem whose source label is \texttt{conic} in the article (source0.tex lines 198--203, source label \texttt{conic}), delivered together with the article's own complete original proof of it (source0.tex lines 204--316, one \texttt{proof} environment of 113 lines). No mathematics is written, repaired or re-derived: statement and proof are the article's own text line for line. Required context retained uncounted: none. Every definition, display and label of those lines is retained unchanged. Omitted from this package and not redistributed: 1--197, 317--587. Nothing inside a retained excerpt is omitted or altered: every retained body is delivered exactly as the article's lines are written, so no omission receipt of any kind accompanies this package. Citations: the retained excerpts carry no citation command at all, so no bibliography entry of the article is transcribed and no reference is redistributed. Reference adaptation: none was needed and none was made; every reference inside the delivered unit resolves inside it, so no unresolved reference or citation is produced. Printed slips kept literal: no printed word, symbol, hypothesis or number of the counted statement or of its proof was altered. Layout and class: the article's own class is \texttt{article} with packages including geometry, color, latexsym, amsmath, amssymb, amsthm, graphicx, subfigure, overpic; the wrapper is the lane's pinned \texttt{amsart} editorial wrapper at 10pt on A4, which restates the theorem-like environments the retained text needs and adds the article's own macro definitions that the retained text needs (\texttt{\textbackslash RR}, \texttt{\textbackslash p}), with extra wrapper packages (bm, bbm, dsfont, xspace, cancel, mathtools). The delivered numbering is the unit's own. Assets: no retained span contains a figure environment, a graphics-inclusion command, a PDF-page inclusion, a table or a TikZ picture; no image file is copied into this package, the package declares no asset dependency and no image file is redistributed with it. Licence: version arXiv:2510.12191v1 of this article is distributed under the Creative Commons Attribution 4.0 International licence, as recorded in the arXiv licence metadata for this exact version and shown on the abstract page https://arxiv.org/abs/2510.12191v1. A full rights notice is delivered at licenses/arxiv-2510.12191-CC-BY-4.0.txt. Characters: every retained excerpt is pure ASCII, so the delivered engine, XeLaTeX with its default Latin Modern text font, reports no missing character.
\begin{lem}\label{conic}
Let $\p=(p_1,p_2,p_3), \p'=(p_1',p_2',p_3')\in  (\RR^2)^3$, and assume that $p_1,p_2,p_3$ are pairwise distinct.
Then either\\
(i)  $\p$ and $\p'$ are congruent, or \\
(ii) $\Sigma_{\p,\p'}$ is one-dimensional, and each of $\sigma_{\p,\p'}$ and $\sigma'_{\p,\p'}$ is an algebraic curve of degree at most two.
\end{lem}

\begin{proof}
By definition, for $(q,q')\in \Sigma_{\p,\p'}$ we have
\begin{align*}
\|p_i-q\|^2=&\|p_i'-q'\|^2,~~~i=1,2,3
\end{align*}
or
$$
\|p_i\|^2-2p_i\cdot q+\|q\|^2=\|p_i'\|^2-2p_i'\cdot q'+\|q'\|^2,~~~i=1,2,3.
$$
Subtracting  the $3$rd equation from each of the first two equations, we get the system
\begin{align}
\|p_1\|^2-\|p_3\|^2-2(p_1-p_3)\cdot q&= \|p'_1\|^2-\|p'_3\|^2-2(p_1'-p_3')\nonumber\\
\|p_2\|^2-\|p_3\|^2-2(p_2-p_3)\cdot q&= \|p'_2\|^2-\|p'_3\|^2-2(p_1'-p_3')\label{sysK3n}\\
\|p_3\|^2-2p_3\cdot q+\|q\|^2&=\|p_3'\|^2-2p_3'\cdot q'+\|q'\|^2.\nonumber
\end{align}
The system can be rewritten as
\begin{align*}
\tfrac12u-A q&=\tfrac12 v-B q', \\
\|p_3\|^2-2p_3\cdot q+\|q\|^2&=\|p_3'\|^2-2p_3'\cdot q'+\|q'\|^2 ,
\end{align*}
where $A$ (resp., $B$) is a $2\times 2$ matrix whose $i$th row equals $p_i-p_3$ 
(resp., $p'_i-p'_3$), for $i=1,2$, and 
\begin{align*}
u & = \left(\begin{matrix}
           \|p_1\|^2-\|p_3\|^2 \\
          \|p_2\|^2-\|p_3\|^2 
         \end{matrix}
         \right) , \quad 
         v   = \left(\begin{matrix}
           \|p'_1\|^2-\|p'_3\|^2 \\
          \|p'_2\|^2-\|p'_3\|^2 
         \end{matrix}
         \right) 
\end{align*}
are vectors in $\RR^2$. 

Assume first that $p_1',p_2',p_3'$ are not collinear. So the matrix $B$ is invertible and we have 
$$
q'=B^{-1}Aq+w,
$$
for $w = \frac12 B^{-1}(v-u) \in\RR^2$. Let $T(q):=B^{-1}Aq+w$. 
So $(q,q')\in \Sigma_{\p,\p'}$ if and only if $q'=T(q)$. Plugging this is the 3rd equation in \eqref{sysK3n} we get
\begin{equation}\label{projq}
\|p_3\|^2-2p_3\cdot q+\|q\|^2=\|p_3'\|^2-2p_3'\cdot T(q)+\|T(q)\|^2,
\end{equation}
which defines $\sigma_{\p,\p'}$. Note that this gives a conic section unless \eqref{projq} is trivial, i.e., the zero equation. 
Note that for this to happen the quadratic part has to be zero. That is, 
$$
\langle q,q\rangle =\langle B^{-1}Aq,B^{-1}A q\rangle 
$$
%or
%$$
%\langle q,q\rangle =\langle q,(B^{-1}A)^{tr}B^{-1}A q\rangle 
%$$
or
$$
\langle q,\left(I-(B^{-1}A)^{tr}B^{-1}A\right)q\rangle =0
$$
This defines a trivial quadratic equation if and only if 
$$
I-(B^{-1}A)^{tr}B^{-1}A=0
$$
or
%$$
%(B^{-1}A)^{tr}B^{-1}A=I
%$$
%or 
$$
(B^{-1}A)^{tr}=(B^{-1}A)^{-1},
$$
%so 
%$$
%\text{$B^{-1}A$ is an orthogonal matrix},
%$$
which means that $T$ is an isometry of $\RR^2$. 
This implies that $\p$ and $\p'$ are congruent.
We conclude that $\sigma_{\p,\p'}$ is either conic section given by  \eqref{projq} or $\p$ and $\p'$ are congruent, which completes the proof for the case that $p_1',p_2',p_3'$ are non-collinear. 

By symmetry, same analysis applies also when $p_1,p_2,p_3$ are non-collinear. 
So we need to prove the lemma for the case that each of the triples $p_1,p_2,p_3$ and $p_1',p_2',p_3'$ is collinear. In this case we may assume without loss of generality that $p_3=p_3'=(0,0)$, and $p_1=(a,0)$, $p_1'=(a',0)$, $p_2=(b,0)$, and $p_2'=(b',0)$. Note that our assumption that $p_1,p_2,p_3$ are pairwise distinct implies $b\neq a\neq 0$. Writing $q=(x,y)$ and $q'=(x',y')$, the system of equations defining $\Sigma_{\p,\p'}$ in this case becomes
\begin{align*}
(x-a)^2+y^2&=(x'-a')^2+(y')^2\\
(x-b)^2+y^2&=(x'-b')^2+(y')^2\\
x^2+y^2&=(x')^2+(y')^2
\end{align*}
or
\begin{align*}
-2ax+a^2&=-2a'x'+(a')^2\\
-2bx+b^2&=-2b'x'+(b')^2\\
x^2+y^2&=(x')^2+(y')^2.
\end{align*}
Recalling our assumption that $p_1,p_2,p_3$ are distinct, we have $b\neq a\neq 0$. Thus the first two equations give 
\begin{align*}
x&=(a'/a)x'-(a')^2/(2a)+a/2\\
x&=(b'/b)x'-(b')^2/(2b)+b/2.
\end{align*}
So we either get unique values for $x,x'$ or otherwise $a'/a=b'/b$. In the former case both $\sigma_{\p,\p'}$ and $\sigma'_{\p,\p'}$  are lines parallel to the $y$-axis, and so in particular (ii) holds. 
In the latter case, write $a'/a=b'/b=t$. Assume first that $t\neq 0$. Then 
\begin{align*}
x&=tx'+(a/2)(1-t^2)\\
x&=tx'+(b/2)(1-t^2),
\end{align*}
and so either $t=\pm 1$, in which case $\p$ and $\p'$ are congruent, or $a=b$ which is a contradiction to our assumption that $p_1,p_2,p_3$ are distinct. 
So we may assume that $t=0$, which means that $a'=b'=0$. 
Assume $(q,q')\in \Sigma_{\p,\p'}$ and write $q=(x,y)$ and $q'=(x',y')$. 
Let $r:=(x')^2+(y')^2$. So $q=(x,y)$ must satisfy the system
\begin{align*}
(x-a)^2+y^2&=r\\
(x-b)^2+y^2&=r\\
x^2+y^2&=r.
\end{align*}
However, recalling that $b\neq a\neq 0$, this system has no solutions. So $\Sigma_{\p,\p'}$ is empty in this case. This completes the proof of the lemma. 
\end{proof}

\end{document}
